\subsection{Phantom Revenue}
When the relaxation solution $\solution = (\fleetAssignVar^*, \shareVar^*, \outsideShareVar^*)$ satisfies both resource availability and inter-departure constraints, it might still be illegal if it relies on "Phantom Revenue." This occurs when the revenue claimed by the aggregated super itineraries cannot be physically realized by any valid combination of their constituent original itineraries in the full discretization network.

Due to the highly coupled nature of the Sales-Based Linear Program (SBLP), the outside option share $\outsideShareVar$ acts as a dynamic variable. A local decrease in the attractiveness of available itineraries can shift demand to the outside option, thereby implicitly raising the choice bounds for other itineraries and reallocating passenger flows. Because of this endogeneity, evaluating Phantom Revenue via static slack analysis is insufficient. Instead, we formulate a Revenue Feasibility Subproblem (RFS) to establish a strict necessary and sufficient condition.

Given the fixed physical service assignment $\fleetAssignVar^*$, we seek to detect whether the target revenue contribution of each super itinerary $R$, defined as $Rev_{Rpl}^* = \demand_{mp} \fare_{Rl} \shareVar_{Rpl}^*$, can be attained by assigning valid shares $\shareVar_{rpl}$ to its constituent original itineraries $r \in \eqClassByPathSig_{(R)}$. We introduce continuous slack variables $s_{Rpl} \ge 0$ to capture the revenue shortfall for each super itinerary. The RFS is formulated as the following linear program:
\full{\begin{align}
RFS(\solution)=\min_{\shareVar, \outsideShareVar, s} \quad & \sum_{m\in \marketSet} \sum_{p\in \psgTypeSet} \sum_{l\in \fareClassSet} \sum_{R \in \itinSuperInMarket_m} s_{Rpl} \label{eq:rfs_obj}\\
\text{s.t.} \quad & \sum_{r \in \eqClassByPathSig_{(R)}} \demand_{mp} \fare_{rl} \shareVar_{rpl} + s_{Rpl} = \demand_{mp} \fare_{Rl} \shareVar_{Rpl}^* \quad \forall R \in \itinSuperInMarket_m, p\in \psgTypeSet, l\in \fareClassSet \label{eq:rfs_target}\\
& \outAttr_{mp} \shareVar_{rpl} \le \attr_{rpl} \outsideShareVar_{mp} \quad \forall m\in \marketSet, p\in \psgTypeSet, l\in \fareClassSet, r \in \itinInMarket_m \label{eq:rfs_ratio}\\
& \left(\frac{\adjOutAttr_{mp}}{\outAttr_{mp}}\right) \outsideShareVar_{mp} + \sum_{r \in \itinInMarket_m} \sum_{l\in \fareClassSet} \left(\frac{\adjAttr_{rpl}}{\attr_{rpl}}\right) \shareVar_{rpl} \le 1 \quad \forall m\in \marketSet, p\in \psgTypeSet \label{eq:rfs_limit}\\
& \sum_{m\in \marketSet}\sum_{p\in \psgTypeSet}\sum_{l\in \fareClassSet} \sum_{r \in \itinUsingArc{m}{a}} \demand_{mp} \shareVar_{rpl} \le \sum_{f\in \fleetSet} \fleetCap_f \fleetAssignVar_{af}^* \quad \forall a \in \flightArc \label{eq:rfs_cap}\\
& \shareVar_{rpl} \ge 0, \quad \outsideShareVar_{mp} \ge 0, \quad s_{Rpl} \ge 0 \nonumber
\end{align}}

Constraint (\ref{eq:rfs_target}) forces the underlying original itineraries to meet the relaxation revenue, using $s_{Rpl}$ to absorb any unavoidable deficit. Constraints (\ref{eq:rfs_ratio})-(\ref{eq:rfs_cap}) ensure that the recovered passenger flow strictly adheres to the physical capacity of $\fleetAssignVar^*$ and the original discrete choice bounds. 

The relaxation solution $\solution$ contains Phantom Revenue (and is therefore illegal) \textit{if and only if} the optimal objective value of the RFS is strictly greater than zero (i.e., $\sum s_{Rpl}^* > 0$). The time slots (i.e. a set of segment-time pairs) corresponding to the super itineraries with positive slack $s_{Rpl}^* > 0$ are the culprits that contribute to the Phantom Revenue. By adding time points to the discretization $\discretization^{2}$ we can ensure that all these phantom revenue violations will not occur in the refined model. Denote the refined discretization as $\discretization^{3}$. The refine procedure will be described in the Subsection~\ref{sec:refine_process}.